\documentclass[a4paper,12pt]{article}

\usepackage[T2A]{fontenc}
\usepackage[utf8]{inputenc}
\usepackage[russian]{babel}
\usepackage{amsmath,amssymb,mathtools}
\usepackage{helvet}
\renewcommand{\familydefault}{\sfdefault}
\usepackage{icomma}
\usepackage[a4paper,margin=19mm,headheight=15pt,footskip=12mm]{geometry}
\usepackage{microtype}
\usepackage{tikz}
\usetikzlibrary{calc,arrows.meta,positioning,shapes.geometric}
\usepackage{tabularx,booktabs,longtable,array}
\usepackage{enumitem}
\usepackage{hyperref}
\usepackage{parskip}
\usepackage{fancyhdr}
\usepackage{setspace}
\usepackage{xcolor}

\definecolor{accent}{HTML}{2A6F6B}
\definecolor{darkaccent}{HTML}{184B48}
\definecolor{textgray}{HTML}{2D3338}
\definecolor{softgray}{HTML}{F2F4F3}
\definecolor{rulegray}{HTML}{C9D0CE}

\hypersetup{hidelinks}
\geometry{includeheadfoot}
\setlength{\parindent}{0pt}
\setlength{\parskip}{0.55em}
\onehalfspacing
\setlist[itemize]{leftmargin=1.45em,itemsep=0.25em,topsep=0.3em}
\setlist[enumerate]{leftmargin=1.6em,itemsep=0.35em,topsep=0.3em}
\renewcommand{\arraystretch}{1.32}

\pagestyle{fancy}
\fancyhf{}
\fancyhead[L]{\small\color{textgray}Анна Баннова}
\fancyhead[R]{\small\color{textgray}Степени и преобразования}
\fancyfoot[C]{\small\color{textgray}\thepage}
\renewcommand{\headrulewidth}{0.3pt}
\renewcommand{\footrulewidth}{0pt}

\newcommand{\topic}{Степени: приведение выражений к общему основанию}
\newcommand{\lessondate}{04.10.26}
\newcommand{\student}{Анна Баннова}
\newcommand{\teacher}{Лёвин Артём Александрович}

\newcommand{\sectionline}[1]{%
  \vspace{0.65em}%
  {\large\bfseries\color{darkaccent}#1}\par
  \vspace{0.18em}\color{rulegray}\hrule\color{textgray}\vspace{0.45em}%
}

\newcommand{\formula}[1]{%
  \begin{center}
    \color{darkaccent}\large$#1$
  \end{center}
}

\tikzset{
  flow/start/.style={ellipse,draw=accent,line width=0.9pt,fill=softgray,align=center,minimum width=40mm,minimum height=11mm,font=\small},
  flow/action/.style={rectangle,rounded corners=2pt,draw=accent,line width=0.9pt,align=center,text width=112mm,minimum height=13mm,font=\small,inner sep=5pt},
  flow/decision/.style={diamond,aspect=2.25,draw=accent,line width=0.9pt,align=center,text width=52mm,minimum height=15mm,font=\small,inner sep=2pt},
  flow/arrow/.style={-{Stealth[length=3mm,width=2mm]},line width=0.85pt,draw=darkaccent},
  flow/label/.style={font=\small\bfseries,fill=white,inner sep=1pt,text=darkaccent}
}

\begin{document}
\color{textgray}

% -------------------- TITLE --------------------
\thispagestyle{empty}
\begin{center}
  \vspace*{22mm}
  {\small\color{accent}\bfseries ЧЕК-ЛИСТ}\par
  \vspace{7mm}
  {\Huge\bfseries\color{darkaccent}Степени}\par
  \vspace{4mm}
  {\Large\bfseries Приведение выражений к общему основанию}\par
  \vspace{13mm}
  {\large \lessondate}\par
  \vfill
  {\large\bfseries \teacher}\par
  \vspace{2mm}
  {\large эксклюзивно для Анны Банновой}\par
  \vspace{16mm}
  \begin{tikzpicture}[x=1cm,y=1cm]
    \draw[accent,line width=1.15pt]
      (0,0) .. controls (3.2,1.1) and (7.4,-1.0) .. (10.8,0.15)
            .. controls (12.4,0.7) and (14.1,0.55) .. (15.6,0.2);
    \fill[accent] (15.6,0.2) circle (1.4pt);
  \end{tikzpicture}
  \vspace*{8mm}
\end{center}
\clearpage

% -------------------- CORE CHECKLIST --------------------
\sectionline{1. Главная идея}

Большинство выражений со степенями становится заметно проще, если привести участвующие числа к \textbf{одинаковым основаниям}, а затем последовательно применить свойства степеней. Рабочий ориентир: после каждого шага выражение должно становиться компактнее.

\sectionline{2. Формулы, которые нужно узнавать с первого взгляда}

Для положительных оснований и действительных показателей:

\begin{align*}
 a^m\cdot a^n &= a^{m+n},\\
 \frac{a^m}{a^n} &= a^{m-n},\\
 \left(a^m\right)^n &= a^{mn},\\
 (ab)^n &= a^n b^n,\\
 a^{-n} &= \frac{1}{a^n},\\
 \sqrt[k]{a} &= a^{1/k}.
\end{align*}

\textbf{Важно.} В задачах занятия основания положительные. Это позволяет безопасно работать с дробными и иррациональными показателями по указанным формулам.

\sectionline{3. Что переводить в степенной вид}

\begin{itemize}
  \item \textbf{Составное основание:} $6=2\cdot3$, $8=2^3$, $42=6\cdot7$.
  \item \textbf{Корень:} $\sqrt[12]{2}=2^{1/12}$.
  \item \textbf{Десятичную дробь:} $0{,}5=\frac12=2^{-1}$, \quad $0{,}2=\frac15=5^{-1}$.
  \item \textbf{Обычное число при необходимости:} $5=5^1$.
\end{itemize}

Разложение выбирайте \textbf{под задачу}. Если рядом уже стоят основания $6$ и $7$, то для $42$ часто выгоднее записать $42=6\cdot7$, а не раскладывать до простых множителей.

\sectionline{4. Скобки: памятка «ВУПС»}

Если внутри действия стоит \textbf{целое выражение}, сохраните его границы скобками:

\begin{center}
\begin{tabular}{@{}ll@{}}
\toprule
\textbf{В} & вычитаемое выражение \(\rightarrow\) взять в скобки;\\
\textbf{У} & умножаемое выражение \(\rightarrow\) взять в скобки;\\
\textbf{П} & подставляемое выражение \(\rightarrow\) взять в скобки;\\
\textbf{С} & скобки сохраняют знаки и структуру.\\
\bottomrule
\end{tabular}
\end{center}

Например,
\[
  a^{\sqrt7+2}:a^{\sqrt7-1}
  =a^{(\sqrt7+2)-(\sqrt7-1)}.
\]
Скобки у вычитаемого показателя защищают знак «минус» перед всем выражением.

\sectionline{5. Базовый пример: составное основание}

Упростим:
\[
\frac{2^{7/2}\cdot3^{11/2}}{6^{9/2}}.
\]

\begin{align*}
\frac{2^{7/2}\cdot3^{11/2}}{6^{9/2}}
&=\frac{2^{7/2}\cdot3^{11/2}}{(2\cdot3)^{9/2}}\\
&=\frac{2^{7/2}\cdot3^{11/2}}{2^{9/2}\cdot3^{9/2}}\\
&=2^{7/2-9/2}\cdot3^{11/2-9/2}\\
&=2^{-1}\cdot3=\frac32.
\end{align*}

\textbf{Самопроверка:} сначала было несколько оснований и дробь, в финале осталось одно обычное число.

\clearpage

\sectionline{6. Корни и дробные показатели}

Пример:
\[
\left(\frac{2^{1/3}\cdot2^{1/4}}{\sqrt[12]{2}}\right)^2.
\]

Сначала переводим корень в степень, затем собираем одинаковые основания:
\begin{align*}
\left(\frac{2^{1/3}\cdot2^{1/4}}{2^{1/12}}\right)^2
&=\left(2^{\frac13+\frac14-\frac1{12}}\right)^2\\
&=\left(2^{1/2}\right)^2=2.
\end{align*}

При сложении дробных показателей приводите их к общему знаменателю обычным способом.

\sectionline{7. Отрицательный показатель — это промежуточный результат}

\formula{a^{-n}=\dfrac{1}{a^n}}

Например,
\[
2^{-2}=\frac{1}{2^2}=\frac14=0{,}25.
\]

Если в ходе преобразований появился отрицательный показатель, это обычно означает, что выражение уже почти упрощено. Перепишите его без отрицательной степени в финале, если этого требует формат ответа.

\sectionline{8. Иррациональные показатели не мешают формулам}

Показатель вида $\sqrt5+2$ остаётся обычным показателем степени. Если основания одинаковы, показатели складываются и вычитаются так же:
\[
\frac{3^{\sqrt5+2}\cdot3^{3-\sqrt5}}{3^4}
=3^{(\sqrt5+2)+(3-\sqrt5)-4}=3.
\]

Корни в показателях часто \textbf{сокращаются алгебраически}; переводить сам показатель $\sqrt5$ в какой-либо другой вид не требуется.

\sectionline{9. Ищите короткий маршрут}

Рассмотрим:
\[
\frac{6^{\sqrt3}\cdot7^{\sqrt3}}{42^{\sqrt3-1}}.
\]

Здесь особенно выгодно заметить общий показатель в числителе:
\begin{align*}
6^{\sqrt3}\cdot7^{\sqrt3}
&=(6\cdot7)^{\sqrt3}=42^{\sqrt3},\\
\frac{42^{\sqrt3}}{42^{\sqrt3-1}}
&=42^{\sqrt3-(\sqrt3-1)}=42.
\end{align*}

\textbf{Стратегия:} если формула сразу уменьшает количество объектов в выражении, применяйте её. Полное разложение на простые множители требуется не всегда.

\sectionline{10. Типичные ошибки}

\begin{itemize}
  \item Складывать показатели при \textbf{разных основаниях}: $2^m\cdot5^n$ нельзя заменить одной степенью без дополнительного преобразования.
  \item Потерять скобки при вычитании составного показателя: $m-(n-1)\ne m-n-1$.
  \item Забыть умножить показатели в $(a^m)^n$.
  \item Неверно раскрыть $(ab)^n$: степень относится к \textbf{каждому} множителю.
  \item Считать отрицательный показатель ошибкой вместо перехода к обратному числу.
  \item Раскладывать основание слишком глубоко, когда уже видна более короткая общая база.
\end{itemize}

\clearpage

% -------------------- FLOWCHART --------------------
\thispagestyle{plain}
\begin{center}
{\Large\bfseries\color{darkaccent}Алгоритм: как упрощать выражения со степенями}
\end{center}
\vspace{7mm}

\begin{center}
\begin{tikzpicture}[node distance=6.5mm]
  \node[flow/start] (start) {Начало};
  \node[flow/action,below=of start,text width=108mm] (scan) {Просмотрите выражение: отметьте одинаковые основания, составные числа, корни, десятичные дроби, произведения и частные.};
  \node[flow/action,below=of scan,text width=108mm] (convert) {Выберите выгодное представление. Приведите основания к общему виду: например, $8=2^3$, $\sqrt[12]{2}=2^{1/12}$, $0{,}5=2^{-1}$, $42=6\cdot7$, если рядом уже есть $6$ и $7$.};
  \node[flow/action,below=of convert,text width=108mm] (open) {Раскройте степенные скобки: $(ab)^n=a^nb^n$, $(a^m)^n=a^{mn}$. Если сразу виден короткий ход через общий показатель, используйте его.};
  \node[flow/action,below=of open,text width=108mm] (combine) {Соберите одинаковые основания: при умножении показатели сложите, при делении — вычтите. Составное вычитаемое заключите в скобки.};
  \node[flow/action,below=of combine,text width=108mm] (finish) {Сократите показатель. Отрицательную степень при необходимости замените обратным числом.};
  \node[flow/decision,below=8mm of finish,text width=48mm] (simple) {Остался очевидный шаг упрощения?};
  \node[flow/start,below=8mm of simple] (end) {Финальный ответ};

  \draw[flow/arrow] (start) -- (scan);
  \draw[flow/arrow] (scan) -- (convert);
  \draw[flow/arrow] (convert) -- (open);
  \draw[flow/arrow] (open) -- (combine);
  \draw[flow/arrow] (combine) -- (finish);
  \draw[flow/arrow] (finish) -- (simple);
  \draw[flow/arrow] (simple) -- node[flow/label,right] {нет} (end);
  \draw[flow/arrow] (simple.west) -- ++(-24mm,0) |- node[flow/label,pos=0.22,left] {да} (scan.west);
\end{tikzpicture}
\end{center}

\clearpage

% -------------------- TRAINING --------------------
\sectionline{Тренировочный блок — Путь А}

Упростите выражения. Записывайте ключевые преобразования: к какому основанию приводите и какую формулу применяете.

\begin{enumerate}
  \item $2^{3/4}\cdot2^{5/4}$.
  \item $\dfrac{5^{7/3}}{5^{4/3}}$.
  \item $\left(3^{2/5}\right)^5$.
  \item $4^{3/2}$.
  \item $\left(\dfrac{2^{1/3}\cdot2^{1/4}}{\sqrt[12]{2}}\right)^2$.
  \item $\dfrac{2^{7/2}\cdot3^{3/2}}{6^{5/2}}$.
  \item $\dfrac{3^{\sqrt5+2}\cdot3^{3-\sqrt5}}{3^4}$.
  \item $\dfrac{8^{2-\sqrt7}\cdot2^{3\sqrt7-1}}{2^2}$.
  \item $\dfrac{(0{,}5)^{\sqrt{10}-1}}{2^{-\sqrt{10}}}$.
  \item $\dfrac{6^{\sqrt3}\cdot7^{\sqrt3}}{42^{\sqrt3-1}}$.
\end{enumerate}

\vfill
\textbf{Контроль перед сдачей:}
\begin{itemize}
  \item одинаковы ли основания перед сложением или вычитанием показателей;
  \item раскрыты ли степенные скобки;
  \item не потеряны ли скобки у составного вычитаемого;
  \item переведён ли отрицательный показатель в удобную финальную форму;
  \item действительно ли последний шаг упрощает выражение.
\end{itemize}

\clearpage

% -------------------- ANSWERS --------------------
\sectionline{Ответы}

\begin{center}
\begin{tabularx}{0.76\linewidth}{@{}>{\bfseries}c X >{\raggedleft\arraybackslash}p{32mm}@{}}
\toprule
№ & Тип основного приёма & Ответ \\
\midrule
1 & умножение степеней с одинаковым основанием & $4$ \\
2 & деление степеней с одинаковым основанием & $5$ \\
3 & степень степени & $9$ \\
4 & приведение основания $4=2^2$ & $8$ \\
5 & корень $\rightarrow$ дробный показатель & $2$ \\
6 & разложение $6=2\cdot3$ & $\dfrac23$ \\
7 & иррациональные показатели, одинаковое основание & $3$ \\
8 & приведение $8=2^3$ и раскрытие степени & $8$ \\
9 & $0{,}5=2^{-1}$ & $2$ \\
10 & общий показатель и $6\cdot7=42$ & $42$ \\
\bottomrule
\end{tabularx}
\end{center}

\vfill
\begin{center}
{\small\color{textgray}\teacher\ эксклюзивно для Анны Банновой · \lessondate}
\end{center}

\end{document}
